\chapter{Introducción}
\label{ch:introduccion}

\section{Propósito del documento}

Este documento especifica el marco LLM-Sec, un instrumento de evaluación de la postura de ciberseguridad para plataformas que integran modelos de lenguaje de gran escala (LLM) en arquitecturas de generación aumentada por recuperación (RAG). El marco combina un modelo de madurez de seis niveles con una taxonomía de 62 controles operacionales, proporcionando al evaluador criterios cuantificables y reproducibles para determinar el estado de seguridad de un sistema RAG.

\section{Alcance}

El marco LLM-Sec está diseñado para evaluar sistemas RAG que cumplan las siguientes condiciones:

\begin{itemize}[leftmargin=2em]
  \item Utilizan un modelo de lenguaje como componente de inferencia.
  \item Incorporan una base de datos vectorial como fuente de contexto.
  \item Operan en un entorno donde la entrada del usuario se procesa en lenguaje natural.
  \item El modelo opera sin agencia sobre el sistema que lo aloja: no invoca funciones, no accede al sistema de archivos, no establece conexiones de red ni ejecuta código. Su única salida es texto generado.
\end{itemize}

Esta última condición delimita la arquitectura evaluable a un flujo RAG conversacional y excluye las arquitecturas agénticas, en que el modelo dispone de herramientas que le permiten actuar sobre su entorno. Tales sistemas introducen vectores de amenaza adicionales, como la ejecución de acciones no autorizadas derivadas de una instrucción inyectada, que requieren controles de aislamiento y validación de invocación no contemplados en la taxonomía de este marco.


El marco cubre cuatro categorías de amenaza del OWASP LLM Top 10:2025 pertinentes al flujo RAG: LLM01 (Prompt Injection), LLM02 (Sensitive Information Disclosure), LLM07 (System Prompt Leakage) y LLM08 (Vector and Embedding Weaknesses).

\section{Audiencia}

\begin{itemize}[leftmargin=2em]
  \item \textbf{Evaluadores de seguridad:} profesionales que realizan auditorías de sistemas RAG.
  \item \textbf{Equipos de desarrollo:} responsables de implementar controles en el ciclo de vida del sistema.
  \item \textbf{Oficiales de cumplimiento:} encargados de verificar la adherencia a normativa vigente.
  \item \textbf{Dirección:} tomadores de decisiones que requieren una visión ejecutiva del nivel de riesgo.
\end{itemize}

\section{Referencias normativas}

El marco LLM-Sec se ancla a las siguientes referencias:

\begin{infobox}[title=Referencias principales]
  \begin{tabularx}{\textwidth}{@{}l X@{}}
    \textbf{OWASP}              & LLM Top 10:2025 \\
    \textbf{ISO/IEC 33001}      & Evaluación de capacidad de procesos \\
    \textbf{ISO/IEC 27002:2022} & Controles de seguridad de la información \\
    \textbf{NIST AI RMF}        & Marco de gestión de riesgos de IA \\
    \textbf{Ley 21.663}         & Marco de Ciberseguridad (Chile) \\
    \textbf{Ley 21.719}         & Protección de Datos Personales (Chile) \\
    \textbf{Boletines 15.869-19 y 16.821-19} & Texto refundido, Proyecto de Ley de Inteligencia Artificial (Chile) \cite{leyiaproyecto} \\
  \end{tabularx}
\end{infobox}

\section{Convenciones del documento}

A lo largo de este documento se utilizan las siguientes convenciones visuales:

\begin{itemize}[leftmargin=2em]
  \item Los controles se identifican con \texttt{C01}\,--\,\texttt{C62}.
  \item Los niveles de madurez se representan con etiquetas: \badgeNivel{0} \badgeNivel{1} \badgeNivel{2} \badgeNivel{3} \badgeNivel{4} \badgeNivel{5}.
  \item Los tipos funcionales se distinguen por color: \badgePreventivo{} \badgeDetectivo{} \badgeCorrectivo{} \badgeGobernanza{}.
  \item Las fuentes de evidencia se indican como: \badgeConfig{} \badgeExperimental{} \badgeDeclarativo{}.
\end{itemize}

Los términos técnicos empleados en este documento se definen en el Apéndice~\ref{app:glosario}.